% Sidebar sobre panic:
% panic é o último recurso do kernel: o impossível aconteceu e
% o kernel não sabe como prosseguir. No xv6, panic faz...

\chapter{Interrupções e drivers de dispositivo}
\label{CH:INTERRUPT}

Um
\indextext{driver}
é o código em um sistema operacional que gerencia um determinado dispositivo:
ele configura o hardware do dispositivo,
instrui o dispositivo a realizar operações,
lida com as interrupções resultantes
e interage com processos que possam estar esperando
por E/S do dispositivo.
O código de drivers pode ser complicado
porque um driver executa concorrentemente com o dispositivo que gerencia.
Além disso, o driver deve entender a interface de hardware do dispositivo,
que pode ser complexa e mal documentada.

Dispositivos que precisam de atenção do sistema operacional geralmente podem ser
configurados para gerar interrupções, que são um tipo de trap.
O código de tratamento de traps do kernel reconhece quando um dispositivo
gerou uma interrupção e chama o manipulador de interrupção do driver.

Muitos drivers de dispositivo executam código em dois contextos: uma \indextext{metade superior}
que roda na thread do kernel de um processo, e uma \indextext{metade inferior}
que executa no momento da interrupção. A metade superior
é chamada por meio de chamadas de sistema como {\tt read} e {\tt write}, que querem
que o dispositivo realize operações de E/S. Esse código pode solicitar ao hardware que inicie uma
operação (por exemplo, pedir ao disco que leia um bloco); então o código espera pela
conclusão da operação. Eventualmente, o dispositivo conclui a
operação e gera uma interrupção. O manipulador de interrupção do driver,
atuando como a metade inferior,
descobre qual operação foi concluída, acorda um processo que estava esperando,
se for o caso, e instrui o hardware a iniciar o trabalho
da próxima operação pendente.

\section{Código: Entrada do console}

O driver de console \fileref{kernel/console.c}
é uma ilustração simples da estrutura de um driver. O
driver de console aceita caracteres digitados por um humano, via o hardware da porta serial \indextext{UART}
conectado ao RISC-V. O driver de console acumula uma
linha de entrada por vez, processando caracteres especiais como
backspace e control-u. Processos de usuário, como o shell, usam
a chamada de sistema {\tt read} para obter linhas de entrada do console.
Quando você digita algo no xv6 executando no QEMU, suas teclas são entregues ao
xv6 por meio do hardware UART simulado pelo QEMU.

O hardware UART com o qual o driver interage é um chip 16550
~\cite{ns16550a} emulado pelo QEMU. Em um computador real, um 16550
gerenciaria um link serial RS232 conectado a um terminal ou outro
computador. Quando executado no QEMU, ele está conectado ao seu teclado e
display.

O hardware UART é exposto ao software como um conjunto de registradores de controle \indextext{mapeados em memória}.
Ou seja, há alguns endereços físicos que
o hardware RISC-V conecta ao dispositivo UART, de forma que operações de leitura e escrita
interajam com o hardware do dispositivo em vez da RAM.
Os endereços mapeados em memória para o UART começam em 0x10000000, ou {\tt UART0}.
Há alguns poucos registradores de controle do UART, cada um com largura
de um byte. Seus deslocamentos a partir de {\tt UART0} são definidos em constantes.
O registrador {\tt LSR} contém bits que indicam se há caracteres de entrada
aguardando leitura pelo software. Esses
caracteres (se houver) podem ser lidos a partir do
registrador {\tt RHR}. Cada vez que se lê um caractere, o hardware do UART
o remove de uma fila interna (FIFO) de caracteres esperando,
e limpa o bit de "pronto" em {\tt LSR} quando a FIFO estiver vazia.
O hardware de transmissão do UART é largamente independente do hardware de recepção;
se o software escrever um byte em {\tt THR},
o UART transmitirá esse byte.

A função {\tt main} do xv6 chama {\tt consoleinit}
para inicializar o hardware. Esse código configura o UART para gerar
uma interrupção de recepção cada vez que um byte de entrada for recebido,
e uma interrupção de \indextext{transmissão completa} cada vez que a transmissão for finalizada.

O shell do xv6 lê do console por meio de um descritor de arquivo.
As chamadas de sistema {\tt read} percorrem o kernel até chegar em {\tt
  consoleread}, que espera pela chegada de entrada (via interrupções) e seu
armazenamento em {\tt cons.buf}, copia a entrada para o espaço do usuário e (após
uma linha completa ter sido recebida) retorna ao processo de usuário. Se o usuário
ainda não tiver digitado uma linha completa, qualquer processo leitor esperará na
chamada {\tt sleep}
(o Capítulo~\ref{CH:SCHED} explica os detalhes da {\tt sleep}).

Quando o usuário digita um caractere, o hardware UART instrui o RISC-V
a gerar uma interrupção, o que ativa
o tratador de traps do xv6.
O tratador de traps chama {\tt devintr},
que examina o registrador {\tt scause} do RISC-V para descobrir que
a interrupção veio de um dispositivo externo.
Em seguida, ele consulta uma unidade de hardware chamada PLIC
\cite{riscv:priv}
para descobrir qual dispositivo gerou a interrupção.
Se foi o UART, {\tt devintr} chama {\tt uartintr}.

A função {\tt uartintr}
lê quaisquer caracteres de entrada aguardando no hardware UART
e os entrega para {\tt consoleintr}.
Ela não espera por caracteres, pois futuras entradas gerarão novas interrupções.
A função {\tt consoleintr} acumula caracteres de entrada em
{\tt cons.buf}
até que uma linha inteira seja recebida.
{\tt consoleintr} trata backspace e alguns outros caracteres
de forma especial.
Quando uma nova linha chega, {\tt consoleintr} acorda um
{\tt consoleread} que esteja esperando (se houver algum).

Uma vez acordado, {\tt consoleread} verá uma linha completa em {\tt
  cons.buf}, a copiará para o espaço do usuário e retornará (via a máquina de chamadas
de sistema) para o espaço do usuário.

\section{Código: Saída do console}

Uma chamada de sistema {\tt write} em um descritor de arquivo conectado ao console
eventualmente chega em
{\tt uartputc}.
O driver de dispositivo mantém um buffer de saída ({\tt uart\_tx\_buf})
para que os processos que escrevem não precisem esperar que o UART termine
de enviar; em vez disso, {\tt uartputc} adiciona cada caractere ao buffer,
chama {\tt uartstart} para iniciar a transmissão pelo dispositivo (caso ele ainda não esteja transmitindo),
e retorna. A única situação em que {\tt uartputc}
espera é se o buffer já estiver cheio.

Cada vez que o UART termina de enviar um byte, ele gera uma interrupção.
{\tt uartintr} chama {\tt uartstart}, que verifica se o dispositivo
realmente terminou de enviar e fornece ao dispositivo o próximo caractere
de saída no buffer. Assim, se um processo escreve múltiplos bytes no
console, normalmente o primeiro byte será enviado pela chamada de
{\tt uartstart} dentro de {\tt uartputc}, e os bytes restantes no buffer serão enviados
por chamadas a {\tt uartstart} feitas por {\tt uartintr} à medida que as
interrupções de transmissão completa chegam.

Um padrão geral a ser notado é o desacoplamento entre a atividade do dispositivo
e a atividade dos processos, feito por meio de buffers e interrupções. O driver
de console pode processar entrada mesmo quando nenhum processo está esperando
para lê-la; uma leitura subsequente verá a entrada. Da mesma forma, processos
podem enviar saída sem precisar esperar pelo dispositivo. Esse desacoplamento
pode melhorar o desempenho, permitindo que processos sejam executados
concorrentemente com a E/S de dispositivos, e é particularmente importante
quando o dispositivo é lento (como o UART) ou requer atenção imediata
(como ao exibir caracteres digitados). Essa ideia às vezes é chamada de
\indextext{concorrência de E/S}.

\section{Concorrência em drivers}

Você pode ter notado chamadas a {\tt acquire} em {\tt consoleread}
e em {\tt consoleintr}. Essas chamadas adquirem um lock, que protege
as estruturas de dados do driver de console contra acesso concorrente.
Há três perigos de concorrência aqui: dois processos em
CPUs diferentes podem chamar {\tt consoleread} ao mesmo tempo;
o hardware pode solicitar a uma CPU que entregue uma interrupção de console (na verdade
do UART) enquanto essa CPU já está executando dentro de
{\tt consoleread};
e o hardware pode entregar uma interrupção de console em
uma CPU diferente enquanto {\tt consoleread} está em execução.
O Capítulo~\ref{CH:LOCK} explica como usar locks
para garantir que esses perigos não resultem em comportamentos incorretos.

Outra forma em que a concorrência exige cuidado em drivers é que um
processo pode estar esperando por entrada de um dispositivo, mas a interrupção
que sinaliza a chegada dessa entrada pode ocorrer quando outro processo (ou
nenhum processo) está em execução. Portanto, manipuladores de interrupção não devem
presumir nada sobre o processo ou código que interromperam. Por
exemplo, um manipulador de interrupção não pode chamar com segurança {\tt copyout}
com a tabela de páginas do processo atual. Manipuladores de interrupção
tipicamente fazem apenas um trabalho mínimo (por exemplo, apenas copiar os dados
de entrada para um buffer) e acordam o código da metade superior para
fazer o restante.

\section{Interrupções de temporizador}

O xv6 usa interrupções de temporizador para manter sua noção do
tempo atual e para alternar entre processos que utilizam intensamente a CPU.
As interrupções de temporizador vêm de um hardware de clock conectado a cada CPU RISC-V.
O xv6 programa o hardware de clock de cada CPU para interromper periodicamente.

O código em {\tt start.c}
permite o acesso em modo supervisor aos registradores de controle do temporizador,
e então solicita a primeira interrupção de temporizador.
O registrador de controle {\tt time} contém um contador que o
hardware incrementa a uma taxa constante; isso serve como uma
noção de tempo atual. O registrador {\tt stimecmp}
contém o tempo no qual a CPU deve gerar uma interrupção de temporizador;
definir {\tt stimecmp} para o valor atual de {\tt time} mais {\it x}
agendará uma interrupção para {\it x} unidades de tempo no futuro.
Para a emulação RISC-V do {\tt qemu}, 1000000 unidades de tempo equivalem
a aproximadamente um décimo de segundo.

As interrupções de temporizador chegam por meio de {\tt usertrap} ou {\tt kerneltrap}
e {\tt devintr}, como outras interrupções de dispositivo.
Interrupções de temporizador chegam com os bits baixos do {\tt scause}
configurados como cinco; {\tt devintr} em {\tt trap.c} detecta essa situação
e chama {\tt clockintr}.
Essa última função incrementa {\tt ticks},
permitindo que o kernel acompanhe a
passagem do tempo. O incremento ocorre em apenas uma CPU, para evitar que o tempo
avance mais rápido se houver múltiplas CPUs.
{\tt clockintr} acorda qualquer processo que esteja esperando na chamada de sistema {\tt sleep},
e agenda a próxima interrupção de temporizador escrevendo em
{\tt stimecmp}.

{\tt devintr} retorna 2 para uma interrupção de temporizador
com o objetivo de indicar a {\tt kerneltrap}
ou {\tt usertrap} que devem chamar {\tt yield}, de forma que
as CPUs possam ser multiplexadas entre processos prontos para executar.

O fato de que código do kernel pode ser interrompido por uma interrupção de temporizador que
força uma troca de contexto via {\tt yield} é parte da razão pela qual
o código inicial em {\tt usertrap} é cuidadoso ao salvar o estado, como {\tt
  sepc}, antes de habilitar interrupções. Essas trocas de contexto também significam
que o código do kernel deve ser escrito com a consciência de que ele pode ser movido
de uma CPU para outra sem aviso prévio.

\section{Mundo real}

O xv6, como muitos sistemas operacionais, permite interrupções e até trocas de contexto
(via {\tt yield}) enquanto executa no kernel. A razão para isso é manter tempos
de resposta rápidos durante chamadas de sistema complexas que executam por muito tempo.
No entanto, como mencionado acima, permitir interrupções no kernel é a fonte de
alguma complexidade; como resultado, alguns sistemas operacionais permitem
interrupções apenas durante a execução de código em modo usuário.

Dar suporte a todos os dispositivos de um computador típico em toda sua complexidade
é muito trabalho, pois há muitos dispositivos, os dispositivos têm muitos
recursos, e o protocolo entre o dispositivo e o driver pode ser complexo
e mal documentado. Em muitos sistemas operacionais, os drivers representam
mais código do que o próprio núcleo do kernel.

O driver de UART recupera dados um byte por vez lendo os registradores de controle do UART;
esse padrão é chamado de \indextext{E/S programada}, já que
o software é quem controla a movimentação dos dados. A E/S programada é simples,
mas lenta demais para ser usada em altas taxas de dados. Dispositivos que precisam
mover grandes quantidades de dados em alta velocidade geralmente usam
\indextext{acesso direto à memória (DMA)}.
O hardware do dispositivo DMA grava diretamente os dados recebidos na RAM e lê
os dados a serem enviados diretamente da RAM. Discos e dispositivos de rede modernos usam DMA.
Um driver para um dispositivo DMA prepara os dados na RAM e, então, usa uma única
escrita em um registrador de controle para instruir o dispositivo a processar os dados preparados.

Interrupções fazem sentido quando um dispositivo precisa de atenção em momentos
imprevisíveis e não com muita frequência. No entanto, interrupções possuem alto custo
de CPU. Assim, dispositivos de alta velocidade, como controladoras de disco e rede,
usam técnicas para reduzir a necessidade de interrupções. Uma técnica é gerar
uma única interrupção para um lote inteiro de requisições de entrada ou saída.
Outra técnica é o driver desabilitar completamente as interrupções e verificar
periodicamente se o dispositivo precisa de atenção. Essa técnica
é chamada de \indextext{polling}. Polling faz sentido se o dispositivo opera
a uma taxa alta, mas desperdiça tempo de CPU se o dispositivo estiver
na maior parte do tempo ocioso. Alguns drivers alternam dinamicamente entre
polling e interrupções dependendo da carga atual do dispositivo.

O driver de UART copia os dados recebidos primeiro para um buffer no kernel
e depois para o espaço do usuário. Isso faz sentido em baixas taxas de dados,
mas essa cópia dupla pode reduzir significativamente o desempenho para dispositivos
que geram ou consomem dados muito rapidamente. Alguns sistemas operacionais
conseguem mover dados diretamente entre buffers em espaço de usuário e o
hardware do dispositivo, frequentemente com DMA.

Como mencionado no Capítulo~\ref{CH:UNIX}, o console aparece para
as aplicações como um arquivo comum, e aplicações leem entrada e escrevem
saída usando as chamadas de sistema \lstinline{read} e \lstinline{write}.
Aplicações podem querer controlar aspectos de um dispositivo que não podem ser
expressos pelas chamadas padrão do sistema de arquivos (por exemplo,
ativar/desativar o buffer de linha no driver de console).
Sistemas Unix oferecem a chamada de sistema \lstinline{ioctl} para esses
casos.

Alguns usos de computadores exigem que o sistema responda em um
tempo limitado. Por exemplo, em sistemas críticos para segurança, perder um
prazo pode levar a desastres. O xv6 não é adequado para ambientes
de tempo real rígido. Sistemas operacionais para tempo real rígido tendem a ser
bibliotecas que se ligam à aplicação de forma que permita uma
análise para determinar o pior caso de tempo de resposta.
O xv6 também não é adequado para aplicações de tempo real suave,
quando ocasionalmente perder um prazo é aceitável, pois o escalonador do xv6
é muito simplista e há caminhos de código no kernel em que as interrupções
ficam desabilitadas por muito tempo.

\section{Exercícios}

\begin{enumerate}

\item Modifique {\tt uart.c} para não usar interrupções. Pode ser necessário
modificar {\tt console.c} também.

\item Adicione um driver para uma placa Ethernet.

\end{enumerate}
